\section{清华岁月：打破信息差与评价体系}

\subsection{开源作业——"比捐楼更有用"}

2016年入学清华后，翁家翌做的第一件"名留青史"的事情是：\textbf{把自己所有的作业和收集到的上古学长资料全部在 GitHub 上开源}。

\begin{importantbox}{信息平权的理念}
翁家翌认为，信息差在清华是一个很有用的生存工具，但每个人都应该平等地拥有这些信息。很多有能力的人不擅长搜集资料，如果能给他们信息平权的机会，他们在清华会活得更好。开源作业不是为了让人抄，而是让后来者不必花十几二十个小时钻牛角尖——把时间省下来做更有价值的事。
\end{importantbox}

这个 GitHub repo 在清华计算机系广为流传。翁家翌半开玩笑地说："你随便抓个计算机系的学弟问，他可能不认识捐信息楼的人，但他认识翁家翌——毕竟大家都是看我的作业活的。"他还自嘲道："（在清华的知名度）比捐楼有用。"

这件事也引发了争议。一些学长学姐反对开源作业，认为这会让后来的学生不再独立思考。但翁家翌坚持认为这是对的——他的目标不是让人抄作业，而是让有能力但不擅长搜集资料的人不再疲于奔命，能用更多时间做更想做的事。不然的话，学生经常花十几二十个小时钻牛角尖，又不敢问助教，对学习的收益非常低。

\subsection{科研启蒙：误打误撞进入强化学习}

大二时，翁家翌开始科研，选择了朱军老师的实验室。当时有三个方向可选：贝叶斯方法、对抗生成网络（GAN）、强化学习（RL）。他本来想做图像相关的 GAN，但因为不知道哪个选项对应 GAN，误选了强化学习。

\begin{warningbox}{科研选择中的随机性}
翁家翌坦言自己选择 RL 方向完全是"random"的——不知道 GAN 是第几个选项就随便选了。但他后来发现 RL 是"打游戏的东西"，觉得挺有意思就一直做了下去。许多看似改变人生的决定，在当时可能只是一个随机事件。
\end{warningbox}

他的第一个 RL 项目是用神经网络通关一个90年代的游戏 VizDoom，并拿了冠军。但他\textbf{并不享受}这个过程——环境太单一，需要疯狂 overfit，调参难度比 CV 高出十到一百倍，大部分情况下算法都不能 work。

这段经历让他得出两个关键认知：
\begin{enumerate}
    \item 当时 RL 研究的核心瓶颈不在算法创新，而在 heuristic 的调参
    \item 他更擅长也更喜欢做\textbf{支撑研究的基础设施}，而非研究本身
\end{enumerate}

翁家翌描述了当时 RL 研究的痛苦：你必须用一些非常 heuristic 的方法去避免各种 corner case。改算法其实没有那么本质——对于人类来说很简单的 task（比如识别障碍物），对 AI 来说是完全不同的难度。他甚至有一种"生理上的排斥反应"——不是不擅长调参，而是\textbf{不喜欢}做这件事。这种清晰的自我认知，让他把重心从"如何做好研究"转向了"如何让研究更顺利"——也就是后来的 infra 之路。

\subsection{三个兴趣方向与评价体系的觉醒}

除了 AI，翁家翌在清华期间还对\textbf{图形学}和\textbf{网络安全}感兴趣。

在图形学方面，他最初的兴趣来自初中看过的科幻电影《创战纪》（Tron），被其电影特效深深震撼——"如果有一天我能做出这种特效，或者构建自己的虚拟世界，那就圆满了"。他在图形学课上拿了全班唯二的 A+，发明了一个减少迭代收敛次数的新算法，渲染了一张当时前所未有的 16K 分辨率图像——在他之前完全没有人渲过 16K 的图。他觉得图形学是"对现实世界的 hacking"，可以以自己的视角构建脑中想象的场景。

在网络安全方面，他觉得"hacker 非常酷"，业余时间搞了很多相关的东西，给学校修了不少校园网的 bug。比如他和一个学长发现了一个漏洞，可以一分钱甚至不要钱下载成绩单（原本每次十块钱），下载几次后把这个 bug 反馈给了学校教务部门。

但最终他选择专注 AI 方向——"如果你要搞科研的话，那还是专心比较好，不能脚踏两条船"。

\begin{knowledgebox}{朱军老师的三指标评价体系}
翁家翌的导师朱军提出了一个不同于 GPA 的评价体系：计算机系学生的价值可以用三个指标衡量——\textbf{论文、比赛成绩、GitHub 三位数以上的 star}。这个评价体系深刻影响了翁家翌，让他意识到可以在开源社区创造与众不同的价值。
\end{knowledgebox}

在 GPA 问题上，翁家翌的策略是\textbf{最低限度投入}——计算好当前分数，确保达到自己设定的标准（比如87分的 B+），多一分都不愿意花时间。他认为 GPA 是一个三四年后就不用写在简历上的东西，不值得投入过多精力。

\subsection{MILA 暑期研究：与 Yoshua Bengio 的交集}

大三暑假，翁家翌通过导师联系到 Yoshua Bengio（2019年图灵奖得主），前往 MILA 做暑期研究。任务是做一个类似 MoE（Mixture of Experts）的东西——有一个 router 选择不同的 path，应用在 language model 上。

这段经历的结果并不理想：当时没有足够的算力和工程能力来 scale up，几块卡根本搞不出来。但回过头看，这次经历让翁家翌同时接触了 RL 和 NLP（Transformer），为后来在 OpenAI 做 RL for LLM 的工作奠定了双重基础。

翁家翌描述了当时的困境：他被一个做 RL 的人派去做 NLP，需要花很长时间入门 Transformer 和 NLP。搓出来的东西没有好效果——现在回过头看，这个东西要 work，需要算力和很强的工程能力来 scale up，一个人几块卡根本不行。"哪怕方向是对的，你也是搞不出来的。"这段暑研经历也没有产出论文，对申请造成了不小的影响。

\begin{warningbox}{先见之明的局限性}
翁家翌强调"马后炮是没有用的"。虽然事后看他在 OpenAI 之前就同时具备了 RL 和 NLP 的经验，但在当时他完全预见不到未来。NLP 在他看来 task 太分散，RL + Transformer 会崩，用几块卡做 MoE 即使方向对了也做不出来。\textbf{很多改变世界的认知，需要足够的 compute 和 engineering 才能被验证。}
\end{warningbox}

\subsection{本章小结}

清华四年，翁家翌完成了三个关键转变：（1）从竞赛思维转向开源和信息平权理念；（2）确立了自己的评价体系——论文、比赛、开源项目，而非 GPA；（3）意识到自己更擅长也更享受做 Infra 而非 Research。这些认知为他后来的职业道路指明了方向。


